\chapter{Limitations and Future Work}
\label{ch:limitations}

\section{Domain Specificity}

Results are measured on competitive Pok\'emon. Transfer to other partially
observed simultaneous-action domains requires re-tokenization and new belief
target definitions. The methods are domain-agnostic but empirical claims are
benchmark-specific.

\section{Play Strength Versus Scientific Metrics}

This thesis does not evaluate live win rate against human players. Offline
action match and belief accuracy answer RQ1--RQ3 as framed. They do not certify
ladder competitiveness. A committee should not read the absence of live wins as
a failed thesis. The intended contribution is measurement of belief learning and
combinatorial generalization under offline imitation.

\section{Label and Reconstruction Noise}

Metamon first-person reconstruction may contain errors on edge-case mechanics.
We mitigate with manual spot checks, corrupted-battle filtering, and comparison
of random-bot versus human statistics.

\section{Belief Labels After Reveal}

Belief metrics focus on pre-reveal turns with masked slots. Including post-reveal
slots would inflate accuracy trivially. Item, ability, and tera accuracies are
trained but not broken out as primary tables because move slots dominate the
label mass.

\section{Compute and Model Scale}

We intentionally use modest model size for Master's-scale hardware. Scaling
width, depth, and context length may narrow gaps to large black-box sequence
policies \citep{grigsby2025metamon} but is left for future work.

\section{Offline RL Stability}

CQL and AMAGO fine-tuning can be unstable on sparse or noisy data. Stage B was
not run. Strong Stage A belief and BC results, including the negative
belief-to-policy transfer finding, suffice for the minimum viable bar in
\Cref{ch:experiments}.

\section{Gen~9 Team Preview}

In Gen~9 OU, all opponent species are visible after team preview. Species-based
belief and Split~B keyed only by species composition may understate harder
combinatorial shift involving unrevealed items, abilities, and move sets.
Split~C and finer set-combination held-out evaluation remain open
(\Cref{ch:results}).

\section{Future Directions}

\begin{itemize}
  \item Live evaluation on \texttt{poke-env} against \texttt{RandomPlayer},
    frozen BC bots, and optionally ladder humans, reported as play-strength
    appendix metrics rather than as the primary thesis claim.
  \item Split~C held-out item, ability, and move-set tuples, plus per-attribute
    belief tables and calibration plots.
  \item Sweeps over $\lambda_\mathrm{bel}$ and larger models to test whether
    belief-to-policy transfer appears at higher capacity.
  \item Conservative offline RL (CQL or AMAGO) initialized from scaled BC
    checkpoints.
  \item Context-window and flat-sequence ablations for RQ4 completeness.
  \item Decision-tree and attention case studies on move-belief errors.
  \item Cross-format evaluation (other generations or tiers) for broader
    generalization claims.
  \item Release of combinatorial split manifests as a standalone benchmark
    contribution.
\end{itemize}
